\begin{theorem}[Эрдёш–Секереш, 1935]
При \(m,n\ge2\)
\[
R(m,n)\le R(m-1,n)+R(m,n-1).
\]
В частности, все числа \(R(m,n)\) конечны (индукция по \(m+n\) с базой \(R(1,n)=R(m,1)=1\)).
\end{theorem}
\begin{proof}
Пусть \(G\) — граф на \(N=R(m-1,n)+R(m,n-1)\) вершинах, \(v\in V\). Положим
\[
A=N(v),\qquad B=V\setminus(N(v)\cup\{v\}),\qquad |A|+|B|=N-1 .
\]
Значит, \(|A|\ge R(m-1,n)\) или \(|B|\ge R(m,n-1)\) (иначе \(|A|+|B|\le N-2\)).
\begin{itemize}\tightlist
\item \(|A|\ge R(m-1,n)\): в \(G[A]\) есть \(K_{m-1}\) — добавив \(v\) (смежна со всеми
  вершинами \(A\)), получаем \(K_m\); либо есть независимое множество размера \(n\).
\item \(|B|\ge R(m,n-1)\): в \(G[B]\) есть \(K_m\); либо независимое множество размера
  \(n-1\), и вместе с \(v\) (не смежна ни с одной вершиной \(B\)) — размера \(n\).
\end{itemize}
По лемме из билета~\ref{ticket:7} этого достаточно.
\end{proof}
\begin{graphillustration}{Разбиение в доказательстве Эрдёша–Секереша: $v$ смежна со всеми вершинами $A$ и ни с одной вершиной $B$. Клика $K_{m-1}$ в $A$ дополняется вершиной $v$ до $K_m$; независимое множество размера $n-1$ в $B$ — до размера $n$.}{fig:08-split}
\begin{tikzpicture}[font=\small]
\node[gi selected] (v) at (0,0) {$v$};
\node[draw,ellipse,minimum width=2.8cm,minimum height=1.8cm] (A) at (-2.8,-1.8) {};
\node[draw,ellipse,minimum width=2.8cm,minimum height=1.8cm] (B) at (2.8,-1.8) {};
\node[gi note] at (-2.8,-3.0) {$A=N(v)$, $|A|\ge R(m-1,n)$?};
\node[gi note] at (2.8,-3.0) {$B=V\setminus N[v]$, $|B|\ge R(m,n-1)$?};
\foreach \x/\y [count=\i] in {-3.6/-1.5,-2.8/-1.2,-2.0/-1.5,-3.2/-2.2,-2.4/-2.2}
  {\node[gi small] (a\i) at (\x,\y) {}; \draw[gi c1] (v) -- (a\i);}
\foreach \x/\y [count=\i] in {2.0/-1.5,2.8/-1.2,3.6/-1.5,2.4/-2.2,3.2/-2.2}
  {\node[gi small] (b\i) at (\x,\y) {}; \draw[gi c2,black!45] (v) -- (b\i);}
\node[gi note,anchor=west] at (5.0,0) {сплошные — рёбра $G$\\штриховые — рёбра $\overline G$};
\end{tikzpicture}
\end{graphillustration}


\begin{theorem}[Усиление при чётных слагаемых]
Если \(R(m-1,n)\) и \(R(m,n-1)\) оба чётны, то \(R(m,n)\le R(m-1,n)+R(m,n-1)-1\).
\end{theorem}
\begin{proof}
Пусть \(N=R(m-1,n)+R(m,n-1)-1\) (нечётно) и в \(G\) на \(N\) вершинах нет ни \(K_m\),
ни независимого множества размера \(n\). Как и выше, для каждой вершины
\(|A|\le R(m-1,n)-1\) и \(|B|\le R(m,n-1)-1\), а \(|A|+|B|=N-1\) — значит, обе оценки
точные и \(d(v)=|A|=R(m-1,n)-1\) нечётно для всех \(v\). Нечётное число вершин
нечётной степени противоречит лемме о рукопожатиях.
\end{proof}
Пример: \(R(3,3)\le R(2,3)+R(3,2)=6\); \(R(3,4)\le R(2,4)+R(3,3)-1=4+6-1=9\)
(обе величины чётны) — точное значение.

\begin{cor}[Биномиальная оценка]
\[
R(m,n)\le\binom{m+n-2}{n-1}=\binom{m+n-2}{m-1}.
\]
\end{cor}
\begin{proof}
Индукция по \(m+n\). Если \(m=1\) или \(n=1\): \(R=1\le\binom{m+n-2}{n-1}\)
(при \(m=1\) биномиальный коэффициент \(\binom{n-1}{n-1}=1\), при \(n=1\) — \(\binom{m-1}{0}=1\)).
При \(m,n\ge2\) по теореме и тождеству Паскаля
\[
R(m,n)\le\binom{m+n-3}{n-1}+\binom{m+n-3}{n-2}=\binom{m+n-2}{n-1}.\qedhere
\]
\end{proof}
\begin{cor}
\(R(n,n)\le\binom{2n-2}{n-1}\sim\dfrac{4^{n-1}}{\sqrt{\pi n}}=O\!\left(\dfrac{4^n}{\sqrt n}\right)\).
\end{cor}
\begin{proof}
Формула Стирлинга: \(\binom{2k}{k}\sim\frac{4^k}{\sqrt{\pi k}}\) при \(k=n-1\).
\end{proof}
